\chapter*{Introduction Générale}
\mtcaddchapter % <-- Indique à minitoc qu'un chapitre fantôme est passé par là
Dans un environnement économique de plus en plus concurrentiel, les entreprises sont amenées à prendre des décisions rapides et pertinentes afin de maintenir leur compétitivité. La multiplication des données générées par les systèmes d'information offre aujourd'hui de nouvelles opportunités pour améliorer les processus décisionnels. Toutefois, la simple disponibilité des données ne suffit pas ; celles-ci doivent être collectées, intégrées, analysées et transformées en informations exploitables pour apporter une réelle valeur ajoutée à l'entreprise.

La Business Intelligence et les technologies de la donnée jouent un rôle essentiel dans cette transformation. Elles permettent de consolider les données provenant de différentes sources, de produire des indicateurs de performance fiables et d'offrir aux décideurs une vision globale de l'activité. Parallèlement, l'émergence des techniques de Machine Learning ouvre de nouvelles perspectives en permettant non seulement d'analyser les situations passées, mais également d'anticiper les événements futurs à travers des modèles prédictifs.

Par la suite, la gestion des stocks constitue un enjeu majeur pour les entreprises de distribution dont une mauvaise anticipation de la demande peut entraîner des ruptures de stock, des pertes financières, une baisse de la satisfaction des clients ou encore une augmentation des coûts liés au surstockage. Il devient alors indispensable de disposer d'outils capables d'assister les responsables dans leurs décisions en fournissant des analyses fiables et des prévisions pertinentes.

C'est dans cette optique que s'inscrit ce projet réalisé au sein de l'entreprise BIOSPHERE où l'objectif principal consiste à concevoir et développer une solution décisionnelle intelligente permettant de centraliser les données de l'entreprise, de suivre les indicateurs clés de performance et d'anticiper les risques liés à la gestion des stocks grâce à l'intégration de techniques de Machine Learning.

Pour atteindre cet objectif, une architecture décisionnelle complète a été mise en place en combinant les approches de Data Engineering, de Business Intelligence et de Data Science. Ce projet s'appuie sur la méthodologie Agile Scrum pour la gestion et l'organisation des développements ainsi que sur la méthodologie CRISP-DM pour la conduite des travaux liés à la modélisation prédictive.

Le présent mémoire est structuré en cinq chapitres :

\begin{enumerate}
    \item Le premier chapitre présente le contexte général du projet, l'entreprise d'accueil ainsi que les concepts théoriques relatifs aux systèmes décisionnels, à la Business Intelligence et au Machine Learning.
    \item Le deuxième chapitre est consacré à l'étude des besoins, à l'analyse fonctionnelle du système ainsi qu'à la planification du projet selon la méthodologie Scrum.
    \item Le troisième chapitre décrit la conception et la mise en œuvre de l'architecture décisionnelle, notamment le processus ETL, le Data Warehouse et les tableaux de bord analytiques.
    \item Le quatrième chapitre présente la phase de modélisation prédictive, depuis la préparation des données jusqu'à l'évaluation des modèles de classification et de régression utilisés pour la prévision des stocks.
    \item Le cinquième chapitre détaille l'intégration des différents composants développés au sein d'une application décisionnelle intelligente ainsi que le déploiement de la solution conformément à la phase Deployment de la méthodologie CRISP-DM.
\end{enumerate}
